% Cuerpo editorial abreviado: preguntas, modelo y planteamiento escalar.
\section{Pregunta, contribución y literatura}\label{sec:compact_question}
Estudiamos cómo una intervención en el precio de las reservas bancarias
afecta la asignación de capital entre firmas heterogéneas mediante
liquidación interbancaria, cartera de los bancos y préstamos garantizados.
Una reducción de la remuneración de reservas no equivale a una caída
mecánica de la tasa de crédito. Tampoco más crédito implica más
inversión: recursos propios, precio del capital y valor pignorable
responden simultáneamente a las condiciones monetarias.

La economía es neokeynesiana, cerrada, con un hogar trabajador, dos
tipos empresariales, bancos comerciales con pagos idiosincráticos,
mercado OTC y banco central. La empresa H es tecnológicamente más
productiva que L, pero enfrenta menor pignorabilidad contractual.
Distinguimos la rama S, donde H tiene garantía activa y L es
\emph{prestataria no restringida} en la fecha inicial, de la rama B,
donde ambas empresas están restringidas. Nuestro experimento
perturba únicamente la \emph{remuneración de reservas}.

La contribución precisa es una caracterización condicional de
complementariedad escalar, una contabilidad explícita de transmisión
bancaria y resultados numéricos reproducibles que no confunden
crédito corriente, capital instalado mañana, producción hoy y
pérdida asignativa. No afirmamos haber demostrado dos equilibrios
competitivos de horizonte infinito ni efectos permanentes de
crecimiento.

La interrelación entre política monetaria, heterogeneidad y mala
asignación no es inédita: véase \citet{gonzalez2026}. Los mercados
interbancarios y pagos dialogan con \citet{bianchi2022} y
\citet{bigioSannikov2023}; las distorsiones de capital con
\citet{midrigan2014} y \citet{gopinath2017}. La novedad
\emph{propuesta}, todavía de alcance condicionado, consiste en
integrar esos canales con una restricción robusta de repago.
El escenario se construye con parámetros ilustrativos, no con
estimación empírica para Perú.

\section{Modelo: tecnología, contrato y liquidez}\label{sec:compact_economy}
Las fechas son \(t\geq 0\), los tipos \(i\in\{H,L\}\) y sus
masas \(\nu_i>0\). Una empresa produce
\begin{equation}\label{eq:compact_technology}
Y_{it}=z_tA_i k_{it}^{\alpha}\ell_{it}^{\beta},
\quad A_H>A_L>0,\quad 0<\alpha,\beta,\quad \alpha+\beta<1.
\end{equation}
El capital \(k_{it}\) se hereda al comienzo de la fecha; la
inversión decidida en \(t\) determina \(k_{i,t+1}\). El minorista
paga \(m_t>0\) por el insumo y elige precios nominales con costos
de ajuste Rotemberg; el trabajo interior satisface
\(m_tF_{\ell,i}=w_t\).

\subsection{Hogar y empresas}
El hogar trabajador ahorra en depósitos y bonos, y maximiza
\[
\mathbb E_0\sum_{t\geq0}\beta_W^t
\left\{u_W(C^W_t)-\chi_WL_t^{1+\varphi}/(1+\varphi)\right\},
\quad u_W'(C)=C^{-\sigma_W}.
\]
Su oferta laboral satisface
\(w_t=\chi_WL_t^\varphi(C_t^W)^{\sigma_W}\).
El empresario tiene utilidad marginal
\(u'_E(c)=c^{-\sigma_E}\) y descuento \(\beta_E\);
no se agrega su consumo al del hogar como si fueran el mismo agente.

Tras producir, vender capital depreciado y pagar la deuda heredada,
el patrimonio disponible de la empresa es
\begin{equation}\label{eq:compact_wealth}
N_{it}=\pi^{op}_{it}+(1-\delta)q_tk_{it}
-\frac{R^L_{t-1}b_{i,t-1}}{\Pi_t}
+\frac{R^D_{t-1}a_{i,t-1}}{\Pi_t}.
\end{equation}
El presupuesto y el contrato exigen
\begin{align}
c_{it}+q_tk_{i,t+1}+a_{it}&=N_{it}+b_{it},
\label{eq:compact_budget}\\
R^L_tb_{it}&\leq H_{it}k_{i,t+1},\label{eq:compact_contract}\\
H_{it}&=\theta_i(1-\delta)
\inf_{s'\in\operatorname{supp}(s_{t+1}\mid s_t)}
\bigl\{\Pi_{t+1}(s')q_{t+1}(s')\bigr\}.\nonumber
\end{align}
\(R_t^L\) es la tasa bruta de repago; \(b_{it}\) es el principal,
\(\theta_i\) es la fracción pignorable del capital \emph{y no}
el coeficiente de encaje. El ínfimo usa el soporte endógeno de
las transiciones admisibles: reemplazarlo por una esperanza requeriría
modelar incumplimiento y recuperación.

Si \(\mu_{it}\geq0\) es el dual del colateral, la condición KKT
de complementariedad es
\begin{equation}\label{eq:compact_kkt}
\mu_{it}(H_{it}k_{i,t+1}-R_t^Lb_{it})=0,\qquad
H_{it}k_{i,t+1}-R_t^Lb_{it}\geq0.
\end{equation}
La igualdad del contrato no implica un dual positivo. En el
subcaso de H con deuda positiva, depósitos empresariales nulos,
garantía activa y \(n_{Ht}:=N_{Ht}-c_{Ht}>0\), la identidad exacta es
\begin{equation}\label{eq:compact_capacity}
k_{H,t+1}=\frac{n_{Ht}}{q_t-H_{Ht}/R_t^L},
\quad b_{Ht}=\frac{H_{Ht}k_{H,t+1}}{R_t^L},
\quad q_t>H_{Ht}/R_t^L.
\end{equation}
A recursos propios, precio del capital y garantía \emph{fijos},
una menor tasa eleva la capacidad financiable; la implicación
para inversión de equilibrio no se deduce de esa derivada parcial.

\subsection{Bancos, pagos, OTC y política monetaria}
Un banco retiene patrimonio \(x\), toma depósitos \(d^D\),
presta \(l^B\) y mantiene reservas previas \(m^B\):
\begin{equation}\label{eq:compact_balance}
l^B+m^B=d^D+x,\qquad
d^D\leq\kappa_Dx,\qquad l^B\leq\kappa_Lx.
\end{equation}
Después del shock de depósitos \(\omega\), el excedente anterior
al OTC es
\begin{equation}\label{eq:compact_surplus}
s^B=m^B-\rho_td^D+
\left(\frac{R_t^D}{R_t^R}-\rho_t\right)\omega d^D.
\end{equation}
Las fracciones de déficit y superávit negociadas dependen de matching
y de la tasa interbancaria; el remanente deficitario acude a la
facilidad del banco central. Su rendimiento de liquidez
\(\chi_t(s^B)\) es lineal por tramos. Bajo precios dados, solvencia
holgada y condiciones interiores, la condición marginal de préstamo es
\begin{equation}\label{eq:compact_bank_foc}
R_t^L-R_t^R=\mathbb E_t[\chi'_t(s^B)]+\nu_{L,t},
\qquad\nu_{L,t}\geq0.
\end{equation}
En quiebres con masa positiva se requieren subgradientes; una
condición marginal necesaria no sustituye la optimalidad dinámica
del banco.

El banco central establece precio y cantidad de reservas con reglas
distintas:
\begin{align}
\log R_t^R&=\varrho_R\log R_{t-1}^R+
(1-\varrho_R)\left[
\log\bar R^R+\phi_\pi\log\Pi_t+
\phi_y\log(Y_t/\bar Y)\right]+\epsilon_t^R,
\label{eq:compact_rate_rule}\\
\log M_t^{\mathrm{pre}}&=\varrho_M\log M_{t-1}^{\mathrm{pre}}
+(1-\varrho_M)\log\bar M+\epsilon_t^M.
\label{eq:compact_quantity_rule}
\end{align}
El coeficiente de encaje \(\rho_t\) y el ancho del corredor
\(\Delta_t=R_t^W-R_t^R\) son instrumentos adicionales. No se
identifica un shock de remuneración con un cambio de encaje.
El ejercicio numérico conserva el encaje en \(0.10\) y altera
solo la innovación de la regla de remuneración:
\begin{equation}\label{eq:compact_shock}
\epsilon_0^R\longmapsto\epsilon_0^R+h,\quad
h=-10^{-5},\qquad
\text{sin nuevas innovaciones agregadas para }t\geq1.
\end{equation}
La convención es de tasa bruta por período ilustrativo, no una
variación anualizada ni una intervención estimada.

\subsection{Cronología y objeto económico}
Se observa el estado inicial y se anuncia la política; hogares,
firmas y bancos eligen consumo, cartera y capital futuro. Luego se
realizan pagos de depósitos, negociación OTC, facilidad y liquidación.
En la siguiente fecha ocurren producción, repago y recomposición
patrimonial. Un equilibrio competitivo infinito requiere factibilidad
por historia, solución de los problemas individuales,
consistencia de precios, expectativas y condiciones terminales.
Una raíz de un sistema de ecuaciones marginales necesario no es
automáticamente un equilibrio competitivo.

\section{Reducción de dos fechas y resultado central}
\label{sec:compact_scalar}
Con estados heredados y anuncios futuros comunes, consideramos dos
ramas de un prefijo no lineal de 33 logvariables,
\[
 Z=\left(\mathcal X_1\in\mathbb R^7,\;
u_0\in\mathbb R^{10},\;
x_2\in\mathbb R^6,\;
u_1\in\mathbb R^{10}\right).
\]
Los sistemas \(F_S(Z,h)=0\) y \(F_B(Z,h)=0\) implementan las 33
condiciones de las dos fechas, con distinta complementariedad
para L al inicio. H está restringida en ambas; L tiene préstamo
positivo y garantía holgada en S, y garantía activa en B. La
continuación local desde \(t=2\) y la regularidad de los
Jacobianos requieren hipótesis y comprobaciones adicionales.

Si \(s\geq0\) representa la holgura normalizada de L, una
eliminación local de las demás ecuaciones induce la función
\(\mathcal M(s,h)\), el valor marginal normalizado de financiación.
El problema escalar es
\begin{equation}\label{eq:compact_complementarity}
0\leq s\ \perp\ \mathcal M(s,h)\geq0.
\end{equation}
El cálculo de fuente aporta como evidencia numérica el signo
\(b=\partial_s\mathcal M(0,0)<0\) y
\(a_R=\partial_h\mathcal M(0,0)<0\).
Son \emph{premisas} de la proposición, no un teorema global de
existencia, selección o unicidad de asignaciones dinámicas.
